\section{Conclusion}
\label{sec:conclusion}

Block-diffusion language models have a staircase cost shape that is
structurally different from the linear cost of autoregressive decoding.
The taxi fare climbs with every token; the subway charges per station.
The shapes cross at $L^{*} = \crossoverL$ output tokens, above which
diffusion is faster. The advantage is large --- fifteenfold at natural
output, eightfold at a thousand tokens --- and in cost terms, that is
a fifteen-to-one gap per million events that does not close with a
bigger model, only with batching.

Against a MoE-matched AR model, this is a fair comparison. Moving
from the dense AR-3B to the MoE AR-30B raised the serial per-token
cost sevenfold, pushing the crossover from 59 to 6 tokens. DG wins at
every practical output length in serial mode.

The advantage is limited where batching is available and where tail
SLAs are tight. But in the batch-1 regime --- serial investigation
chains, autonomous driving, algorithmic trading --- the staircase
shape provides a structural latency advantage that the line cannot
match. The $K$ knob is DG's emergency brake: when the deadline
tightens, DG degrades gracefully --- lower quality, complete canvas.
AR's only response is truncation. Block diffusion is the right choice
when output is long, batch-1 is required, and schema completeness
matters. We have named five conditions and grounded them in concrete
workloads.
